\section{Arquitectura del Sistema y Diagrama}
\label{sec:arquitectura}

El diseño arquitectónico de \textbf{artesa} responde al patrón desacoplado de tres capas (\textit{Three-Tier Architecture}) complementado con una pasarela de acceso inverso (\textit{Reverse Proxy Gateway}) para garantizar alta disponibilidad, aislamiento de responsabilidades y seguridad perimetral.

\subsection{Componentes del Sistema}

\subsubsection{Capa de Presentación (Frontend SPA)}
Desarrollada como una Aplicación de Página Única (\textit{Single Page Application} - SPA) utilizando la biblioteca **React 19** impulsada por el empaquetador ultrarrápido **Vite**. La maquetación y diseño de interfaz se basan en **Tailwind CSS**, complementado con la biblioteca de iconografía accesible **Lucide React**. La navegación del lado del cliente está gestionada por **React Router**, integrando componentes de guardia de ruta (\texttt{ProtectedRoute}) que validan el estado de sesión activa y los permisos de usuario requeridos antes de renderizar cada módulo.

\subsubsection{Capa de Pasarela y Enrutamiento (Gateway)}
En entorno de producción contenerizado, un servidor web ligero **Nginx** actúa como pasarela perimetral en el puerto de exposición estándar (\texttt{:8080}). Nginx cumple dos funciones esenciales:
\begin{enumerate}
    \item Servir de forma optimizada los activos estáticos del frontend precompilados en modo producción con cabeceras de caché eficientes.
    \item Actuar como proxy inverso transparente para el prefijo de ruta \texttt{/api/*}, canalizando las peticiones HTTP hacia el microservicio backend en el puerto interno \texttt{:3000}, evitando problemas de Política de Mismo Origen (\textit{CORS}) y ocultando la topología de red interna.
\end{enumerate}

\subsubsection{Capa de Lógica de Aplicación (Backend API REST)}
Construida sobre el entorno de ejecución **Node.js 22 (LTS)** utilizando el framework web **Express** en modo módulos nativos ECMAScript (ESM). La arquitectura del código interno sigue una separación estricta en cuatro niveles por cada dominio de negocio:
\begin{itemize}
    \item \textbf{Routes:} Define las rutas del protocolo HTTP, asocia los middlewares de seguridad (\texttt{authenticate}, \texttt{requirePermiso}) y valida el esquema de entrada.
    \item \textbf{Controller:} Extrae parámetros de la petición HTTP, invoca al servicio correspondiente y devuelve una respuesta estructurada bajo el estándar \textbf{JSend} (estados \texttt{success}, \texttt{fail} o \texttt{error}).
    \item \textbf{Service:} Orquesta la lógica del negocio, gestiona el contexto de solicitud y coordina la ejecución transaccional.
    \item \textbf{Repository:} Interactúa directamente con la base de datos a través de modelos de **Sequelize** para operaciones CRUD o mediante consultas parametrizadas directas para vistas y funciones de PostgreSQL.
\end{itemize}

\subsubsection{Capa de Persistencia y Motor Transaccional (Database)}
Implementada en el motor de base de datos relacional **PostgreSQL 16**. A diferencia de los enfoques tradicionales que delegan la lógica matemática y transaccional crítica en el código JavaScript del servidor, \textbf{artesa} implementa la regla de oro del proyecto: \textit{la lógica transaccional crítica reside en funciones nativas en PL/pgSQL}. Esta decisión garantiza que las operaciones que involucran dinero, inventario y retenciones impositivas se ejecuten en milisegundos con atomicidad garantizada por el propio motor SQL, inmunes a caídas imprevistas del servidor de aplicaciones.

\subsection{Diagrama de Arquitectura}

La Fig.~\ref{fig:arquitectura} ilustra el flujo de peticiones, la segregación de capas y las fronteras de red entre los componentes de la infraestructura.

\begin{figure}[htbp]
\centering
\resizebox{\columnwidth}{!}{%
\begin{tikzpicture}[
    node distance=1.1cm and 0.8cm,
    auto,
    box/.style={rectangle, draw=black!70, fill=blue!5, thick, rounded corners, text width=2.1cm, align=center, minimum height=1.1cm, font=\scriptsize},
    gateway/.style={rectangle, draw=orange!80, fill=orange!10, thick, rounded corners, text width=2.1cm, align=center, minimum height=1.1cm, font=\scriptsize},
    db/.style={cylinder, cylinder uses custom fill, cylinder end fill=green!20, cylinder body fill=green!5, draw=green!60!black, thick, shape border rotate=90, text width=1.9cm, align=center, minimum height=1.3cm, font=\scriptsize},
    line/.style={draw, -latex', thick}
]

% Nodos
\node [box, fill=purple!10] (client) {\textbf{Cliente Web}\\Navegador SPA};
\node [gateway, right=0.9cm of client] (nginx) {\textbf{Gateway / Nginx}\\Puerto :8080\\Reverse Proxy};
\node [box, right=0.9cm of nginx] (express) {\textbf{Backend API}\\Express (Node 22)\\Puerto :3000};
\node [db, right=0.9cm of express] (postgres) {\textbf{PostgreSQL 16}\\DB: artesa\\PL/pgSQL};

% Rutas del backend
\node [box, fill=yellow!10, below=0.7cm of express, text width=2.8cm] (auth) {\textbf{Middlewares}\\JWT Auth + RBAC\\AsyncLocalStorage};
\node [box, fill=green!10, below=0.7cm of postgres, text width=2.4cm] (sqlfunc) {\textbf{Lógica Crítica}\\Funciones SQL\\Transacciones ACID};

% Conexiones
\path [line] (client) -- node[above, font=\tiny] {HTTP/S} (nginx);
\path [line] (nginx) -- node[above, font=\tiny] {/api/*} (express);
\path [line] (express) -- node[above, font=\tiny] {SQL Pool} (postgres);
\path [line] (express) -- (auth);
\path [line] (postgres) -- (sqlfunc);
\path [line, dashed] (auth) -| node[above, font=\tiny, near start] {Auditoría IP} (postgres);

\end{tikzpicture}%
}
\caption{Diagrama de Arquitectura de Capas de \textbf{artesa}.}
\label{fig:arquitectura}
\end{figure}

\subsection{Mecanismo de Autenticación y Contexto Transaccional}
El sistema implementa autenticación sin estado en el transporte pero revalidada en persistencia:
\begin{enumerate}
    \item Tras autenticarse exitosamente mediante \texttt{POST /api/auth/login}, la API responde emitiendo una cookie de sesión cifrada que transporta el JWT con identificadores únicos (\texttt{jti}, \texttt{id\_usuario}, \texttt{id\_rol}).
    \item Para cada petición subsecuente, el middleware \texttt{authenticate} verifica la validez criptográfica del token y consulta el estado del usuario en la base de datos (\texttt{activo = true}), permitiendo la revocación inmediata de privilegios sin esperar a que el token expire.
    \item Mediante \texttt{AsyncLocalStorage}, se captura la dirección IP y el \textit{User-Agent} del cliente HTTP, inyectándolos en variables de sesión de PostgreSQL (\texttt{app.user\_id}, \texttt{app.client\_ip}) durante la ejecución de \texttt{withTransaction}, asegurando que cada disparador de bitácora relacional conserve el contexto de auditoría completo.
\end{enumerate}
